% from ega2-1-fr.tex
\oldpage[II]{154}
et tels que $\mathcal{I}_0=\mathcal{O}_Y$, le $Y$-schéma
$X=\operatorname{Proj}(\bigoplus_{n\geq0}\mathcal{I}_n)$ est intègre et le
morphisme structural $f:X\to Y$ dominant.
\item[\rm(ii)] Soit $\mathcal{I}$ un Idéal fractionnaire quasi-cohérent de
$\mathcal{R}(Y)$, et soit $X$ le $Y$-schéma éclaté de $Y$ relativement à
$\mathcal{I}$. Si $\mathcal{I}\neq0$, le morphisme structural $f:X\to Y$ est
alors birationnel et surjectif.
\end{enumerate}
\end{proposition}

\begin{enumerate}
\item[\rm(i)] résulte de ce que $\mathcal{S}=\bigoplus_{n\geq0}\mathcal{I}_n$
est une $\mathcal{O}_Y$-Algèbre \emph{intègre}
(\hyperref[II.3.1.12-fr]{3.1.12} et \hyperref[II.3.1.14-fr]{3.1.14}) puisque
pour tout $y\in Y$, $\mathcal{O}_y$ est un anneau intègre
(\hyperref[I.5.1.4-fr]{I, 5.1.4}).
\item[\rm(ii)] D'après (i), $X$ est intègre~; si en outre, $x$ et $y$ sont les
points génériques de $X$ et $Y$, on a $f(x)=y$, et il faut prouver que $k(x)$
est de rang $1$ sur $k(y)$. Or $x$ est aussi le point générique de la fibre
$f^{-1}(y)$~; si $\psi$ est le morphisme canonique $Z\to Y$, où
$Z=\operatorname{Spec}(k(y))$, le préschéma $f^{-1}(y)$ s'identifie à
$\operatorname{Proj}(\mathcal{S}')$, où
$\mathcal{S}'=\psi^*(\mathcal{S})$ (\hyperref[II.3.5.3-fr]{3.5.3}). Mais il est
clair que $\mathcal{S}'=\bigoplus_{n\geq0}(\mathcal{I}_y)^n$, et comme
$\mathcal{I}$ est un Idéal fractionnaire quasi-cohérent non nul de
$\mathcal{R}(Y)$, $\mathcal{I}_y\neq0$
(\hyperref[I.7.3.6-fr]{I, 7.3.6}), d'où $\mathcal{I}_y=k(y)$~;
$\operatorname{Proj}(\mathcal{S}')$ s'identifie donc à
$\operatorname{Spec}(k(y))$ (\hyperref[II.3.1.7-fr]{3.1.7}), d'où la
conclusion.
\end{enumerate}

Nous démontrerons une \emph{réciproque} de
(\hyperref[II.8.1.4-fr]{8.1.4}) dans
(\hyperref[III.2.3.8-fr]{III, 2.3.8}).

\begin{env}[8.1.5]
\label{II.8.1.5-fr}
Revenons à la situation et aux notations de
(\hyperref[II.8.1.1-fr]{8.1.1}). Par définition, les homomorphismes
d'injection $\mathcal{I}_{n+1}\to\mathcal{I}_n$
(\hyperref[II.8.1.1.1-fr]{8.1.1.1}) définissent pour chaque $k\in\mathbf{Z}$ un
homomorphisme injectif de degré $0$ de $\mathcal{S}$-Modules gradués
\begin{equation}
\label{II.8.1.5.1-fr}
  u_k:\mathcal{S}_+(k+1)\to\mathcal{S}(k)~;
\tag{8.1.5.1}
\end{equation}
comme $\mathcal{S}_+(k+1)$ et $\mathcal{S}(k+1)$ sont canoniquement
(TN)-isomorphes, il correspond canoniquement à $u_k$ un homomorphisme injectif
de $\mathcal{O}_X$-Modules (\hyperref[II.3.4.2-fr]{3.4.2})~:
\begin{equation}
\label{II.8.1.5.2-fr}
  \widetilde{u}_k:\mathcal{O}_X(k+1)\to\mathcal{O}_X(k).
\tag{8.1.5.2}
\end{equation}

Rappelons d'autre part (\hyperref[II.3.2.6-fr]{3.2.6}) que l'on a défini des
homomorphismes canoniques
\begin{equation}
\label{II.8.1.5.3-fr}
  \lambda:\mathcal{O}_X(h)\otimes_{\mathcal{O}_X}\mathcal{O}_X(k)
  \to\mathcal{O}_X(h+k)
\tag{8.1.5.3}
\end{equation}
et comme le diagramme
\[
\xymatrix{
  \mathcal{S}(h)\otimes_{\mathcal{S}}\mathcal{S}(k)
    \otimes_{\mathcal{S}}\mathcal{S}(l)\ar[r]\ar[d]
  & \mathcal{S}(h+k)\otimes_{\mathcal{S}}\mathcal{S}(l)\ar[d]\\
  \mathcal{S}(h)\otimes_{\mathcal{S}}\mathcal{S}(k+l)\ar[r]
  & \mathcal{S}(h+k+l)
}
\]
est commutatif, il résulte de la fonctorialité des $\lambda$
(\hyperref[II.3.2.6-fr]{3.2.6}) que les homomorphismes
(\hyperref[II.8.1.5.3-fr]{8.1.5.3}) définissent sur
\begin{equation}
\label{II.8.1.5.4-fr}
  \mathcal{S}_X=\bigoplus_{n\in\mathbf{Z}}\mathcal{O}_X(n)
\tag{8.1.5.4}
\end{equation}

